\documentclass[10pt,a4paper]{amsart}
\usepackage[margin=19mm]{geometry}
\usepackage{amsmath,amssymb,mathtools}
\usepackage[scr=rsfs,cal=euler]{mathalfa}
\usepackage{xfrac}
\usepackage[unicode,hidelinks,bookmarks=false]{hyperref}
\newcommand{\ie}{i.e.\ }
\newcommand{\cf}{cf.\ }
\newcommand{\resp}{respectively\ }
\newcommand{\Subsection}{\subsection}
\providecommand{\nobreakdash}{\nobreak\hbox{-}}
\setlength{\emergencystretch}{2em}
\multlinegap0pt
\usepackage{bm}
\usepackage{bbm}
\usepackage{dsfont}
\usepackage{xspace}
\usepackage{cancel}
\usepackage{mathtools}
\usepackage{amsthm}
\makeatletter
\@ifundefined{theorem}{\newtheorem{theorem}{Theorem}}{}
\@ifundefined{definition}{\newtheorem{definition}[theorem]{Definition}}{}
\makeatother
\title[Planarity criteria for metric graphs]{Planarity criteria for metric graphs}
\author{Alice Brolin, Pavel Kurasov}
\date{Source: arXiv:2601.04050v1 (CC BY 4.0)}
\begin{document}
\maketitle

\noindent\emph{Source and licence.} Planarity criteria for metric graphs. Version arXiv:2601.04050v1 (CC BY 4.0), distributed under the Creative Commons Attribution 4.0 International licence, as recorded in the arXiv licence metadata for this exact version and shown on the abstract page https://arxiv.org/abs/2601.04050v1. Full rights notice: licenses/arxiv-2601.04050-CC-BY-4.0.txt.\par\smallskip

\noindent\textit{Editorial scope.} Counted unit: the Theorem whose source label is \texttt{ThST} in the article (source0.tex lines 695--697, source label \texttt{ThST}), delivered together with the article's own complete original proof of it (source0.tex lines 699--739, one \texttt{proof} environment of 41 lines). No mathematics is written, repaired or re-derived: statement and proof are the article's own text line for line. Required context retained uncounted: eq2 (lines 521--523); eq9 (lines 529--541). Every definition, display and label of those lines is retained unchanged. Omitted from this package and not redistributed: 1--520, 524--528, 542--694, 698--698, 740--1537. Nothing inside a retained excerpt is omitted or altered: every retained body is delivered exactly as the article's lines are written, so no omission receipt of any kind accompanies this package. Citations: the retained excerpts carry no citation command at all, so no bibliography entry of the article is transcribed and no reference is redistributed. Reference adaptation: none was needed and none was made; every reference inside the delivered unit resolves inside it, so no unresolved reference or citation is produced. Printed slips kept literal: no printed word, symbol, hypothesis or number of the counted statement or of its proof was altered. Layout and class: the article's own class is \texttt{amsart} with packages including xcolor, tikz, marginnote, amsmath, amsthm, amsrefs, amssymb, graphicx; the wrapper is the lane's pinned \texttt{amsart} editorial wrapper at 10pt on A4, which restates the theorem-like environments the retained text needs and adds the article's own macro definitions that the retained text needs (none), with extra wrapper packages (bm, bbm, dsfont, xspace, cancel, mathtools). The delivered numbering is the unit's own. Assets: no retained span contains a figure environment, a graphics-inclusion command, a PDF-page inclusion, a table or a TikZ picture; no image file is copied into this package, the package declares no asset dependency and no image file is redistributed with it. Licence: version arXiv:2601.04050v1 of this article is distributed under the Creative Commons Attribution 4.0 International licence, as recorded in the arXiv licence metadata for this exact version and shown on the abstract page https://arxiv.org/abs/2601.04050v1. A full rights notice is delivered at licenses/arxiv-2601.04050-CC-BY-4.0.txt. Characters: every retained excerpt is pure ASCII, so the delivered engine, XeLaTeX with its default Latin Modern text font, reports no missing character.
 \begin{equation} \label{eq2}
\mathbf M (\lambda_2) - {\rm diag}\; \{ \alpha_j \} = - \mathbf A.
\end{equation}

\begin{definition}
A metric graph $ \Gamma $ with delta couplings $ \alpha_j $ at the vertices is called
{\bf admissible} if and only if:
\begin{itemize}
\item the second eigenvalue $ \lambda_2 (\Gamma) $ lies below all singularities 
of the M-function. Since the lowest singularity is $\left( \frac{\pi}{\ell_{\rm max}} \right)^2$ this means that
\begin{equation} \label{eq9}
\lambda_2 (\Gamma) <  \left( \frac{\pi}{\ell_{\rm max}} \right)^2; 
\end{equation}
\item the matrix $ \mathbf Q(\lambda_2) = \mathbf M (\lambda_2 (\Gamma)) - {\rm diag}\, \{\alpha_j \} $
possesses the Strong Arnold Property.
\end{itemize}
\end{definition}

\begin{theorem}  \label{ThST} The Colin de Verdi{\`e}re graph parameters $\mu(G)$ and $ \mu (\Gamma) $ are equal, 
provided $ G $ is the discrete graph associated with the metric graph $ \Gamma$.
\end{theorem}

\begin{proof}
Our goal is to prove that $ \mu(G) = \mu (\Gamma) $. We divide the proof into two steps:
\begin{enumerate}
\item $ \mu (G) \leq \mu (\Gamma) .$
Let $\mathbf{A}$ be a Colin de Verdi{\`e}re matrix with maximal multiplicity of the second eigenvalue. We pick a positive number $\lambda_2$ such that $\sqrt{\lambda_2}<|a_{ij}|$ for all $v^i\sim v^j$ and then we 
construct an operator $L_{\vec{\alpha} \delta}$ with $\lambda_2$ as its second eigenvalue and equation \eqref{eq2} is satisfied. 


Consider any edge, say between two vertices $ v^i $ and $ v^j $. The corresponding entries of the matrices $ - \mathbf Q (\lambda_2) $ and $ \mathbf A $ are given by
$$ -\frac{\sqrt{\lambda_2}}{\sin(\sqrt{\lambda_2} \ell_{ij})} \quad \mbox{and} \quad a_{ij} < 0,$$
respectively.
The function $ \ell \mapsto \frac{\sqrt{\lambda_2}}{\sin(\sqrt{\lambda_2} \ell )}$ with  the domain $\ell \in (0,\frac{\pi}{\sqrt{\lambda_2}})$ has the range $[\sqrt{\lambda_2},\infty)$, and  
hence there is $ \ell_{ij} $ making the entries equal
\begin{equation} \label{equal}
a_{ij} = -\frac{\sqrt{\lambda_2}}{\sin(\sqrt{\lambda_2} \ell_{ij})} ,
\end{equation}
 since we have chosen $ \lambda_2 $ satisfying $\sqrt{\lambda_2}<|a_{ij}|$.
Note that there are two possible values of $ \ell_{ij} $, which can be chosen arbitrarily.


In this way we get a metric graph $ \Gamma $ so that
the matrices $-\mathbf{Q}(\lambda_2)$ and $\mathbf{A}$ have identical entries outside the diagonal. 
To make the two matrices identical $-\mathbf{Q}(\lambda_2)=\mathbf{A}$
we adjust 
the values of  the $\delta$-couplings in $ L_{\vec{\alpha}\delta}$.
Since all singularities of $ \mathbf{Q}$ are to the right of $ \lambda_2$
(condition \eqref{eq9}) and $\xi=0$ is the second lowest eigenvalue of $\mathbf{A}$, $\lambda_2$ is the second eigenvalue of $ L_{\vec{\alpha} \delta}$. $ \mathbf{Q}(\lambda_2)$ possesses the Strong Arnold 
Property since $\mathbf{A}$ does. This accomplishes the proof of the first part.

\item $ \mu (G) \geq \mu(\Gamma) .$ Consider one of the operators $ L_{\vec{\alpha}\delta}$, which maximises 
the multiplicity of $ \lambda_2 $ among admissible metric graphs with standard delta couplings. Then taking
$ \mathbf{A} =  - \mathbf Q(\lambda_2) $
we obtain a Colin de Verdi{\`e}re matrix.
Equation \eqref{eq9}
 implies that $ \sin \sqrt{\lambda}_2 \ell_{ij} >0 $, which in turn implies that non-diagonal elements of $ -\mathbf Q(\lambda_2) $
 are negative for $ v^i \sim v^j$. Moreover $ \xi = 0 $ is the second smallest
eigenvalue of $ \mathbf{A}$. The second assertion follows.

\end{enumerate}

\end{proof}

\end{document}
